\documentclass[11pt,a4paper]{article}

\usepackage[utf8]{inputenc}
\usepackage[T1]{fontenc}
\usepackage[spanish]{babel}
\usepackage[a4paper,margin=2.4cm]{geometry}
\usepackage{amsmath,amssymb}
\usepackage{bm}
\usepackage{booktabs}
\usepackage{enumitem}
\usepackage{xcolor}
\usepackage{hyperref}
\usepackage{tcolorbox}

\hypersetup{colorlinks=true, linkcolor=blue!50!black, urlcolor=blue!60!black}

\newcommand{\BHF}{B_{\mathrm{hf}}}
\newcommand{\dEQ}{\Delta E_Q}
\newcommand{\nuR}{\nu}
\newcommand{\code}[1]{\texttt{\small #1}}

\newtcolorbox{nota}[1][Nota]{colback=blue!4, colframe=blue!45!black,
  fonttitle=\bfseries, title={#1}, left=4pt, right=4pt, top=3pt, bottom=3pt}
\newtcolorbox{aviso}[1][Advertencia]{colback=orange!7, colframe=orange!70!black,
  fonttitle=\bfseries, title={#1}, left=4pt, right=4pt, top=3pt, bottom=3pt}

\title{\textbf{Modelos de relajación magnética en Fitbauer}\\[3pt]
\large Componente fenomenológica y modelo dinámico Blume--Tjon de dos estados}
\author{Departamento de Química Física · UAM}
\date{\today}

\begin{document}
\maketitle

\begin{abstract}\noindent
Este documento describe la primera implementación de relajación magnética en
Fitbauer. Se añaden dos modelos nuevos sin modificar los modelos matemáticos
preexistentes de singlete, doblete, sextete ni distribuciones Hesse--Rübartsch:
(1) una componente fenomenológica \code{Relajacion}, que mezcla un sextete
bloqueado con un doblete superparamagnético aparente, y (2) una componente
\code{BlumeTjon}, que aproxima la relajación dinámica mediante intercambio
simétrico de dos estados $+\BHF \leftrightarrow -\BHF$.
\end{abstract}

\tableofcontents
\bigskip

%======================================================================
\section{Motivación física}

En nanopartículas magnéticas, ferritas u óxidos de hierro parcialmente
superparamagnéticos, el momento magnético puede fluctuar durante la ventana de
observación Mössbauer. Si las fluctuaciones son lentas, el núcleo observa un
campo hiperfino casi estático y el espectro es un sextete. Si son rápidas, el
campo magnético se promedia y aparece un doblete o singlete superparamagnético.
En el régimen intermedio las líneas se ensanchan y coalescen.

La implementación pretende evitar una interpretación automática de tipo
``sextete + doblete = dos fases químicas''. Una fracción tipo doblete puede ser
la manifestación espectral de relajación rápida.

\begin{aviso}
Los modelos descritos aquí son complementarios. Si no se seleccionan las
componentes \code{Relajacion} o \code{BlumeTjon}, los modelos anteriores se
evalúan exactamente por las mismas ramas de código que antes.
\end{aviso}

%======================================================================
\section{Notación común}

La transmisión se escribe, en el límite de absorbente delgado, como
\begin{equation}
  T(v)=b+s\,v-A(v),
\end{equation}
donde $b$ es la base, $s$ la pendiente y $A(v)$ la absorción positiva de la
componente. Para un sextete convencional:
\begin{equation}
  A_{6}(v)=D\sum_{j=1}^{6} I_j L\bigl(v; v_j,\Gamma_j\bigr),
\end{equation}
con profundidad $D$, intensidades $I_j$ y perfil lorentziano
\begin{equation}
  L(v;v_0,\Gamma)=\frac{\Gamma^2}{(v-v_0)^2+\Gamma^2}.
\end{equation}
Las posiciones de primer orden son
\begin{equation}
  v_j = \delta + q_j\,\dEQ + m_j\,\frac{\BHF}{33\,\mathrm{T}},
\end{equation}
donde $m_j$ son las posiciones patrón de $\alpha$-Fe a $33$ T y $q_j$ el patrón
cuadrupolar de primer orden.

%======================================================================
\section{Modelo fenomenológico \code{Relajacion}}

\subsection{Definición}

La componente \code{Relajacion} interpola entre un sextete bloqueado y un
doblete superparamagnético aparente:
\begin{equation}
  A_{\mathrm{rel}}(v)=D\left[
  f_{\mathrm{b,ef}}\,S_6(v)
  +\bigl(1-f_{\mathrm{b,ef}}\bigr)\,D_2^{\ast}(v)
  \right].
  \label{eq:rel-phen}
\end{equation}
$S_6(v)$ es el perfil de sextete por profundidad unitaria y $D_2^{\ast}(v)$ es
el doblete escalado para que su área integrada sea comparable a la del sextete:
\begin{equation}
  D_2^{\ast}(v)=D_2(v)\,
  \frac{\int S_6(v)\,dv}{\int D_2(v)\,dv}.
\end{equation}
Así, $D$ conserva aproximadamente el significado de área/profundidad total.

\subsection{Fracción bloqueada y tasa fenomenológica}

El usuario ve dos parámetros específicos:
\begin{itemize}[leftmargin=*]
  \item $f_{\mathrm{bloq}}$: fracción bloqueada máxima.
  \item $\log_{10}\nu$: tasa de relajación fenomenológica en s$^{-1}$.
\end{itemize}
La fracción que entra realmente en \eqref{eq:rel-phen} es
\begin{equation}
  f_{\mathrm{b,ef}}=f_{\mathrm{bloq}}\,F(\log_{10}\nu),
  \label{eq:fb-eff}
\end{equation}
con
\begin{equation}
  F(x)=\frac{1}{1+\exp\left(\frac{x-x_c}{w}\right)},
  \qquad x_c=8.5,\quad w=0.55.
\end{equation}
Por tanto, tasas lentas ($x\ll x_c$) dejan $F\simeq 1$, mientras que tasas
rápidas ($x\gg x_c$) fuerzan $F\simeq 0$.

Además se aplica un ensanchamiento fenomenológico máximo en el régimen
intermedio:
\begin{equation}
  \Gamma_{\mathrm{ef}}=\Gamma\left[1+1.6\exp\left(-\frac{1}{2}
  \left(\frac{x-x_c}{0.75}\right)^2\right)\right].
\end{equation}

\begin{aviso}[Correlación importante]
$f_{\mathrm{bloq}}$ y $\log_{10}\nu$ no son parámetros completamente
independientes: ambos reducen la parte de sextete observada. En la práctica,
conviene fijar uno de los dos salvo que el espectro tenga información suficiente.
Para estimar fracciones, fijar $\log_{10}\nu$ y liberar $f_{\mathrm{bloq}}$; para
estudiar dinámica, fijar $f_{\mathrm{bloq}}=1$ y liberar $\log_{10}\nu$.
\end{aviso}

%======================================================================
\section{Modelo dinámico \code{BlumeTjon}: dos estados}

\subsection{Matriz de intercambio}

La componente \code{BlumeTjon} implementa una primera versión dinámica con dos
orientaciones opuestas del campo:
\begin{equation}
  +\BHF \leftrightarrow -\BHF,
\end{equation}
con matriz de transición simétrica
\begin{equation}
  W=\begin{pmatrix}-\nu & \nu\\ \nu & -\nu\end{pmatrix}.
\end{equation}
No hay fracción bloqueada explícita: la dinámica está controlada por $\nu$.

\subsection{Perfil de intercambio para una transición}

Para cada línea del sextete se calculan dos frecuencias/velocidades, una para
$+\BHF$ y otra para $-\BHF$, denotadas $v_a$ y $v_b$. Se usa una forma de
intercambio simétrico tipo Bloch--McConnell/Kubo. Definimos
\begin{equation}
  z_a = \Gamma + i(v-v_a),\qquad z_b = \Gamma + i(v-v_b),
\end{equation}
y una tasa de intercambio en unidades de velocidad
\begin{equation}
  k_v = \frac{\nu}{\mathcal{C}_{\nu\to v}},
\end{equation}
donde
\begin{equation}
  \mathcal{C}_{\nu\to v}=\frac{E_\gamma}{h c_{\mathrm{mm/s}}}.
\end{equation}
Aquí $E_\gamma$ es la energía de la transición de $^{57}$Fe y
$c_{\mathrm{mm/s}}$ la velocidad de la luz en mm/s. El perfil de absorción usado
es la parte real de
\begin{equation}
  R(v)=\frac{1}{2}\,\Gamma\,
  \frac{z_a+z_b+4k_v}{z_a z_b+k_v(z_a+z_b)}.
  \label{eq:exchange}
\end{equation}
La contribución total es
\begin{equation}
  A_{\mathrm{BT}}(v)=D\sum_{j=1}^{6} I_j\,\Re\left[R_j(v)\right]_+,
\end{equation}
donde $[\cdot]_+$ recorta pequeñas contribuciones negativas numéricas.

\subsection{Límites}

\begin{itemize}[leftmargin=*]
  \item \textbf{Límite lento}: $k_v\to0$. La ecuación \eqref{eq:exchange}
  tiende a la media de dos lorentzianas estáticas. En el patrón simétrico
  $+\BHF/-\BHF$ reproduce el sextete bloqueado.
  \item \textbf{Intermedio}: $k_v$ comparable a la separación de líneas. Aparece
  ensanchamiento y coalescencia.
  \item \textbf{Rápido}: $k_v$ grande. Las posiciones magnéticas se promedian y
  el espectro tiende a un doblete o singlete según $\dEQ$.
\end{itemize}

%======================================================================
\section{Néel--Arrhenius y distribución de tamaños}

La componente \code{NeelSize} usa el modelo de Néel--Arrhenius para convertir el
volumen de partícula en tasa de relajación. Para partículas esféricas de diámetro
$d$,
\begin{equation}
  V(d)=\frac{\pi d^3}{6},
\end{equation}
y
\begin{equation}
  \tau_N(d,T)=\tau_0\exp\left(\frac{K_{\mathrm{eff}}V(d)}{k_B T}\right),
  \qquad
  \nu(d,T)=\tau_0^{-1}\exp\left(-\frac{K_{\mathrm{eff}}V(d)}{k_B T}\right).
\end{equation}
En la GUI se ajusta $\log_{10}K_{\mathrm{eff}}$ y $\log_{10}\tau_0$ para mantener
rangos numéricos razonables:
\begin{equation}
  \log_{10}\nu_i=-\log_{10}\tau_0
  -\frac{K_{\mathrm{eff}}V(d_i)}{k_BT\ln 10}.
\end{equation}

La distribución de tamaños se modela inicialmente como lognormal. Si $d_{50}$ es
la mediana y $\sigma_d$ la anchura lognormal,
\begin{equation}
  f(d)=\frac{1}{d\sigma_d\sqrt{2\pi}}
  \exp\left[-\frac{(\ln d-\ln d_{50})^2}{2\sigma_d^2}\right].
\end{equation}
El espectro total se discretiza como
\begin{equation}
  A_{\mathrm{Neel}}(v)=D\sum_i w_i\,
  A_{\mathrm{BT}}\bigl(v;\log_{10}\nu_i,T,K_{\mathrm{eff}},d_i\bigr),
\end{equation}
donde $A_{\mathrm{BT}}$ es el perfil de intercambio de dos estados de la sección
anterior. Las partículas grandes tienen barrera $K_{\mathrm{eff}}V$ mayor y por
tanto relajan más despacio; las pequeñas relajan rápido y aportan señal
superparamagnética.

La temperatura de bloqueo aproximada para un diámetro dado se estima con
\begin{equation}
  T_B(d)=\frac{K_{\mathrm{eff}}V(d)}{k_B\ln(\tau_M/\tau_0)},
\end{equation}
siendo $\tau_M$ el tiempo de observación Mössbauer (por defecto se usa
$10^{-8}$ s para el resumen de la GUI).

\begin{aviso}[Identificabilidad]
En un único espectro, $K_{\mathrm{eff}}$, $d_{50}$, $\sigma_d$, $\tau_0$, $\Gamma$
y una posible distribución estática $P(\BHF)$ pueden estar fuertemente
correlacionados. Por eso conviene fijar parámetros con información externa
(TEM, magnetometría, literatura) o usar ajuste global multi-temperatura.
\end{aviso}

%======================================================================
\section{Ajuste global multi-temperatura}

La fase 5 añade un backend puro para ajustar simultáneamente varios espectros a
distintas temperaturas. Para cada espectro $r$ se tiene $T_r(v)$ y temperatura
$T_r$. El modelo comparte los parámetros físicos
\begin{equation}
  \theta_{\mathrm{shared}}=\{\delta,\dEQ,\BHF,\Gamma,K_{\mathrm{eff}},d_{50},
  \sigma_d,\tau_0\},
\end{equation}
y mantiene parámetros locales por espectro, por defecto
\begin{equation}
  \theta_r=\{b_r,s_r,D_r\}.
\end{equation}
La función objetivo es
\begin{equation}
  \chi^2=\sum_r\sum_k
  \left(\frac{T_{r,\mathrm{model}}(v_k;\theta_{\mathrm{shared}},\theta_r,T_r)
  -T_{r,\mathrm{data}}(v_k)}{\sigma_{rk}}\right)^2.
\end{equation}
Esto reduce la degeneración frente a un ajuste de un único espectro: la misma
distribución de tamaños debe explicar simultáneamente el bloqueo a baja
$T$ y el colapso superparamagnético a alta $T$.

%======================================================================
\section{Relación con el ajuste y los informes}

Los parámetros ajustables son los mismos que en un sextete ($\delta$, $\dEQ$,
$\BHF$, anchuras, profundidad e intensidades) más $\log_{10}\nu$ y, sólo en el
modelo fenomenológico, $f_{\mathrm{bloq}}$. En \code{NeelSize} se añaden $T$,
$\log_{10}K_{\mathrm{eff}}$, $d_{50}$, $\sigma_d$, $\log_{10}\tau_0$ y el número
de bins. La GUI muestra también
\begin{equation}
  \tau=\frac{1}{\nu}=10^{-\log_{10}\nu}\;\mathrm{s}
\end{equation}
y, para \code{NeelSize}, una $T_B$ aproximada.

En los informes, la componente \code{Relajacion} debe interpretarse como
``fracción bloqueada / superparamagnética aparente'', no como fases químicas
separadas. La componente \code{BlumeTjon} se reporta como modelo dinámico de dos
estados y \code{NeelSize} como modelo Néel--Arrhenius con distribución de tamaños.

%======================================================================
\section{Validación implementada}

Se añadieron pruebas sintéticas para verificar:
\begin{enumerate}[leftmargin=*]
  \item \code{Relajacion} con $f_{\mathrm{bloq}}=1$ reproduce un sextete.
  \item \code{Relajacion} con $f_{\mathrm{bloq}}=0$ reproduce un doblete escalado
  al área del sextete.
  \item $\log_{10}\nu$ controla los límites lento/rápido del modelo
  fenomenológico.
  \item \code{BlumeTjon} reproduce límites lento e intercambio rápido de forma
  cualitativa.
  \item La distribución lognormal de tamaños se normaliza correctamente y la tasa
  Néel disminuye al aumentar el diámetro.
  \item \code{NeelSize} cambia con la temperatura y el backend global recupera
  un diámetro sintético compartido en un par de espectros.
\end{enumerate}

\begin{nota}[Alcance]
La implementación actual no incluye todavía Néel--Arrhenius, distribución de
tamaños ni ajuste global multi-temperatura. Esas extensiones requieren modelos y
datos adicionales para evitar sobreinterpretar un único espectro.
\end{nota}

\end{document}
